\section{Introduction}
A persistent-memory agent can change internally long before task accuracy reveals anything. Consider two copies of the same interaction trajectory. If one is labeled successful and the other failed before reflection, an outcome-conditioned writer can store different reusable advice although the underlying experience is fixed. That establishes a persistent-state intervention, but not behavioral self-improvement: the different state must still be exposed by retrieval, influence a decision, and survive to an outcome.

We study this gap as \emph{stage-resolved transport}. For a frozen trajectory $\tau$, we switch the released success/failure reflection branch---a bundled writer-protocol intervention that changes reflection instruction together with outcome/reward semantics---and instrument
\[
\text{write }(W)\rightarrow\text{source-item exposure }(E)\rightarrow\text{first-action uptake }(U)\rightarrow\text{outcome }(O).
\]
A separate forced fixed-evidence arm $F$ bypasses native retrieval and asks whether the paired memory content can affect the downstream system when supplied. These surfaces answer different questions and have different units. A write difference proves state change, a retrieval hit proves availability of the source item, and neither proves that the branch-differentiating residual changed the policy. Conversely, a weak endpoint alone cannot tell whether the update failed to write, failed to retrieve, was ignored after retrieval, or changed behavior only sparsely. We call this loss of mechanism information \emph{evaluation coarsening}.

The frozen evidence produces exactly this disagreement across surfaces. Reward-conditioned writing diverges for 20/20 Shopping pairs and 4/4 Reddit pairs; a stronger same-mode wording control leaves a reward-mode excess of 0.105 ($p=0.0078$). Forced supply yields terminal $|\Delta|=0.15625$ ($p=0.00074$). Yet native Shopping retrieves the branch-conditioned source item in 125/172 held-out opportunities while the first-action branch contrast is not supported (TV $=0.06944$, $p=0.5801$, 0/36 modal changes) and terminal transport is sparse ($|\Delta|=0.02083$, $p=0.4289$, 34/36 zero). Reddit again has reliable write divergence but 6/8 zero terminal cells and opposite signs in the two nonzero cells. The strongest supported operational localization is therefore after source-item exposure and before stable first-action uptake. It is an evidence boundary, not a causal attenuation onset or mediation coefficient.

Stage resolution is not merely ``report more metrics.'' The native prerequisite order $O\Rightarrow U\Rightarrow E\Rightarrow W$ permits five binary prefix depths; crossing them with $F\in\{0,1\}$ yields a complete 10-state diagnostic class for this experiment. Exhaustive enumeration of all 32 observed-surface subsets shows that $W/E/U/O/F$ is the unique separating basis: omitting a native coordinate aliases adjacent prefix states, while omitting $F$ aliases the two capacity states at every depth. This paper-specific result explains why coarsening can make genuinely different failure explanations observationally identical.

The diagnostic also creates an engineering question: does the first unsupported measured stage identify a useful repair surface? We test that prediction prospectively rather than assuming it. On 13 outcome-blind pilot states selected before new policy outputs, we replay identical branch memories and upstream packets under native prompting ($A0$), a matched memory-blind decision check ($A1$), and a syntax-matched explicit memory-use check ($A2$). All 312 calls complete with frozen model and packet identities, but mean first-action uptake is identical across arms ($U_{A0}=U_{A1}=U_{A2}=0.0962$), mean $D=U_{A2}-U_{A1}=0$, and only 3/13 states have $D>0$. The preregistered pilot gate therefore stops the corresponding 23-state confirmatory test. The negative result sharpens the engineering lesson: \emph{diagnosis can determine what has earned behavioral credit and what surface merits a probe, without identifying the intervention that will repair it}.

Our contributions are:
\begin{enumerate}
    \item \textbf{Controlled stage-resolved transport.} A same-trajectory success/failure writer intervention follows one induced persistent-state contrast through native exposure, first-action uptake, outcome, and an explicit forced-capacity bypass control.
    \item \textbf{Diagnostic completeness under the experiment's prerequisite class.} The complete 10-state class induced by $O\Rightarrow U\Rightarrow E\Rightarrow W$ and $F$ requires the full $W/E/U/O/F$ observation basis to avoid aliases.
    \item \textbf{An empirical transport boundary.} Write divergence and source-item exposure are established while stable first-action uptake is not, and cross-domain terminal effects are sparse and sign-heterogeneous.
    \item \textbf{A negative diagnosis-guided repair pilot.} A prospectively controlled explicit memory-use check does not selectively increase branch-specific uptake under the frozen pilot gate, separating stage-indexed \emph{credit/diagnosis} from unearned \emph{repair efficacy}.
\end{enumerate}

The resulting claim is deliberately narrower than a new memory-utilization method: \textbf{state divergence is not behavioral authority; availability is not use; diagnosis is not a repair prescription}.
